\section{Introduction}\label{sec:r34introduction}
Dynamic economic decisions depend on valuations of outcomes that have not yet
occurred. A continuation approximation is useful not simply when it predicts
those valuations accurately, but when its errors preserve the decisions and
policy comparisons for which the approximation is used. This distinction is
especially important when the future policy changes, when the controls are
multidimensional, or when recursive preferences amplify continuation errors.
A small fitting loss then need not imply a small economic loss, and a cheap
query need not imply a cheap policy computation.

This paper develops Neural Bellman Operators for policy evaluation and
feasible improvement in controlled economies. A scalar neural continuation
supplies the values and derivatives entering the actor's objective. Evaluation
and improvement remain distinct operations: the continuation evaluates a
specified future, while a feasible current action is chosen against that
continuation. Centering the relevant continuation errors identifies the part
that can change a decision. Full-action Bellman inequalities then propagate
those decision errors to a comparison of entire policies. The original
controlled-diffusion formulation is retained, together with its application
conditions and its economic extensions.

Our constructive result concerns a trainable square-activation continuation
in a recursive capital economy. Economic primitives and a fixed feasibility
witness determine an admissible coefficient envelope, an error allocation,
and a finite hidden-weight training budget before any fitted policy target
is available. Each fit evaluates the future policy actually returned by the
backward construction. The resulting bound compares that policy with all
adapted policies of finite recursive cost, over all real vector actions.
It does not rely on a finite state catalogue or an evaluation of the optimal
policy. A residual-centered evaluator and a separate execution account make
the policy comparison applicable to rounded implementations.

Changing economic futures raise a further question: when can a previously
learned neural continuation be trained again rather than discarded? We prove
a perturbation-stable contraction for warm factors whose Gram matrices need
not commute with the new continuation target. The result admits rectangular
factors, updates of all hidden weights, and a stated error in the whole
implemented update. It supplies a finite cap above an explicit precision
floor. A complementary gradient--curvature bound rules out an uninformative
stationary saddle and translates approximate second-order information into
continuation accuracy. These are properties of specified neural
constructions, not properties inferred from successful optimization traces.

The connection from warm training to policy accuracy requires its own
argument. A warm factor can overestimate as well as underestimate the new
continuation, and errors in evaluating that continuation differ from errors
in fitting it. We construct a joint primitive budget that allows both signs
of approximation and separates Gram error, own-policy evaluation error, and
the actor equation residual. A backward induction establishes the candidate's
coefficient envelope and its full-policy comparison together. The induction
therefore does not assume the successful policy bound it is intended to
prove. With rational primitives and strict moment margins, exact evaluation
and dyadically rounded factor updates give a finite-arithmetic realization.
The economic target remains policy accuracy; admissibility of a warm start
is only one condition for constructing a policy that meets that target.

The numerical studies separate three questions: whether a returned policy
meets its economic tolerance, whether the neural construction can be trained
to that tolerance, and how much complete work the construction requires.
The recursive experiments use continuous Gaussian shocks, vector controls,
and reoptimized futures. Conventional and structural procedures receive the
same available information and are evaluated against the same policy target.
Fitting, failed checks, fallback work, actor solves, verification, and durable
output enter the cost account. The earlier finite-law experiments and their
unfavorable reuse comparisons remain part of the evidence. A constructive
neural guarantee is not used to reverse an already observed cost ranking.

The reliability statements are equally specific. A deterministic training
cap is uniform over its declared admissible factors and update errors. A
small number of fitting seeds does not estimate reliability over an
unspecified population of neural fits. Likewise, exact small-matrix checks
validate the stated arithmetic and policy inequalities; they are not counted
as additional economic observations. The directly specified capital economy
has no discretization-transfer term, whereas transferring its conclusions
to a different continuous model would require that term to be established.

The article proceeds from economic targets and centered decision errors to
full-policy comparison, constructive training, warm reuse, and the numerical
studies. The current technical supplement gives the associated proofs.
The original application manuscript and the earlier article and supplement
are preserved as separately identified companion documents, including their
original results, derivations, and distinct comparison classes.
